% Sections/9_project_management_reflection.tex
\section{Project Management \& Reflection}
\subsection{Project Management}
I structured the project using a Gantt chart and logbook, breaking the work into key milestones. The project was initially scoped as a hardware project in which a drone would be built around the vape dataset. This changed after the further investigation in Section 3, which showed that the project needed to address dataset construction and detection-to-cleanup translation before live UAV hardware could be justified.

\subsection{Personal Reflection}
This project became a genuinely interdisciplinary design-engineering project. It required data engineering, software development, machine learning, dataset construction and human-centred workflow design. The most important lesson was that a detection model alone is not a complete solution. For the work to matter, the output of the model had to connect back to the people responsible for finding, prioritising and removing hazardous litter.

I could have scoped the project much more narrowly as a simple vape detector, but the deeper value came from treating the problem end to end: first understanding the waste and cleanup context, then building the DCE, constructing the UAVape dataset, evaluating the detector, and finally exploring how detections could become practical cleanup tasks. A longer note from the author is included in Appendix~\ref{app:extended_reflection}.
